\subsection{Clinical Pressure--Flow Motivation}
\label{subsec:pressure-flow-motivation}

\paragraph{Anatomy--function distinction.}
Coronary stenosis is the motivating lesion class because a visible narrowing
does not automatically imply a flow-limiting obstruction. Clinical assessment
therefore distinguishes anatomical stenosis detection from functional
pressure--flow assessment
\cite[p.~4401]{DerimayEtAl2021CoronaryStenosisObstruction}
\cite[Chs.~1--2 and~5]{EscanedDavies2017PhysiologicalAssessment}
\cite[Sec.~1]{ShaikhEtAl2025CTFFR}. Coronary lesion-imaging studies can
describe local geometry and lesion length, but those anatomical descriptors do
not by themselves define a pressure--flow output
\cite[Abstract; Methods]{VanDenHoogenEtAl2022CoronaryLesionCTA}. This
anatomy--function split motivates reporting pressure drop and reduced pressure
ratio under declared observation and reference conventions, rather than treating
stenosis severity $s_{\max}$ alone as the modeled functional output.
Section~\ref{sec:resolved-3d-comparison} reports velocity diagnostics only; the
pressure-drop and pressure-ratio conventions fix reduced-output vocabulary for
the model framework.
Figure~\ref{fig:anatomical-functional-motivation} summarizes this motivation
as a schematic distinction between visible narrowing and pressure--flow loss.

Clinical fractional flow reserve (FFR) is an invasive pressure-wire index,
conventionally interpreted in its measurement setting as distal coronary
pressure divided by aortic pressure during maximal hyperemia
\cite[Ch.~5, Sec.~5.5]{EscanedDavies2017PhysiologicalAssessment}
\cite[Abstract]{PijlsEtAl1996FFRFunctionalSeverity}. FAME lesion-level evidence
documented the discordance between angiographic and functional severity and is
the clinical source for the familiar $\mathrm{FFR}\le 0.80$ decision context
\cite[Abstract]{ToninoEtAl2010FAMEAngioFunctional}. Computed tomography-derived
FFR (CT-FFR) studies supply noninvasive comparison vocabulary by inferring
pressure--flow quantities from coronary computed tomography angiography
(CCTA)-derived anatomy and comparing them against invasive FFR or related
reference measurements
\cite[Abstract; Methods]{NorgaardEtAl2014NXTFFRCT}
\cite[Secs.~1.1--1.2]{ShaikhEtAl2025CTFFR}.

\begin{figure}[!htbp]
    \centering
    \captionsetup{aboveskip=2pt,belowskip=2pt}
    \figtikz{anatomical-vs-functional-stenosis}
    \caption{Pressure--flow motivation for reduced pressure outputs. The
    schematic distinguishes an anatomical narrowing from the pressure-loss
    signature used by reduced-model outputs.}
    \label{fig:anatomical-functional-motivation}
\end{figure}

For the selected model gauge-pressure field $p_{\mathrm g}$, define the
observed proximal and distal pressure traces through declared scalar observation
maps:
\begin{flalign*}
    \quad
    P_{\mathrm{prox}}(t)
    \defined
    \mathcal O_{\mathrm{prox}}[p_{\mathrm g}(\cdot,t)],
    \qquad
    P_{\mathrm{dist}}(t)
    \defined
    \mathcal O_{\mathrm{dist}}[p_{\mathrm g}(\cdot,t)]. &&
\end{flalign*}
The pressure-drop waveform is the additive-datum-invariant quantity
\begin{flalign*}
    \quad
    \Delta p(t)
    \defined
    P_{\mathrm{prox}}(t)-P_{\mathrm{dist}}(t). &&
\end{flalign*}
Unlike $\Delta p(t)$, a pressure ratio is not invariant under additive pressure
shifts, so a reduced pressure-ratio output must also declare its pressure
reference, averaging window, observation maps, and denominator admissibility.
Convention~\ref{conv:appendix-pressure-reference-ratio} records those
requirements and the trace-level definition of
$\mathrm{PR}_{1D}(\pi_{\mathrm{PR}})$. Figure~\ref{fig:pressure-tap-convention}
shows the corresponding axial-tap shorthand for an idealized 1D stenosis.

\begin{figure}[!htb]
    \centering
    \captionsetup{aboveskip=2pt,belowskip=2pt}
    \figtikz{stenosis-schematic}
    \caption{Schematic tap and output convention for an idealized stenosis,
    not drawn from a computed case. The taps define $\Delta p(t)$;
        Convention~\ref{conv:appendix-pressure-reference-ratio} supplies the
        pressure-ratio observation maps, averaging window, pressure reference,
        and admissibility rule for $\mathrm{PR}_{1D}(\pi_{\mathrm{PR}})$.}
    \label{fig:pressure-tap-convention}
\end{figure}

\paragraph{Shear diagnostics.}
Pressure drops and pressure ratios summarize lesion-scale pressure--flow loss.
A resolved WSS output belongs to a resolved stress state and must declare the
stress tensor, wall normal, wall location, output operator, and signed tangent
direction when a signed component is reported
\cite[Abstract; Sec.~2, pp.~8--10]{JohnEtAl2017WallShearStressVectorForm}.
Reduced shear proxies are model-tier diagnostics rather than clinical
anatomy--function outputs; Convention~\ref{conv:appendix-shear-diagnostic}
records the declaration requirements for comparisons with traction-based WSS.
